\chapter{Refactoring}
\section{Cosa è?}
Il Refactoring è un metodo strutturato e disciplinato che consiste nel riscrivere o ristrutturare del codice sorgente esistente senza modificarne il comportamento esterno.

Viene eseguito applicando una serie di piccoli passi di trasformazione logica, strettamente combinati con la ripetizione dei test dopo ciascun passo per assicurarsi di non aver "rotto" alcuna funzionalità

\section{Quando applicarlo?}
Di norma, il refactoring andrebbe applicato:
\begin{itemize}
    \item Seguendo la "regola del 3", se una cosa si ripete 3 volte, è il momento di estrarla o rifattorizzarla
    \item Quando si aggiunge una nuova funzionalità al sistema
    \item Quando si corregge un bug
    \item Durante le sessioni di revisione del codice
\end{itemize}

\section{Obiettivi}
L'utilità principale del refactoring è il miglioramento continua della vase di codice e la preparazione al cambiamento futuro.

Gli obiettivi di una buona programmazione prevedono:
\begin{itemize}
    \item La rimozione del codice duplicato
    \item L'abbreviazione di metodi troppo lunghi
    \item Miglioramento generale della chiarezza del sistema
\end{itemize}

\section{Procedura per applicare il refactoring}
Il ciclo di lavoro tipico del refactoring si articola in questi passaggi:
\begin{enumerate}
    \item \textbf{Assicurarsi che i test passino}\\
    Prima di toccare il codice, bisogna avera una suite di test funzionante
    \item \textbf{Individuare un "Code Smell"}\\
    Identificare un sintomo visivo o strutturale che indica la presenza di un problema di progettazione
    \item \textbf{Determinare la miglioria}\\
    Seguire la tecnica di refactoring appropriata
    \item \textbf{Effettuare le migliorie}\\
    Modificare il codice sorgente
    \item \textbf{Eseguire i test}\\
    Verificare che il sistema si comporti ancora in modo corretto
    \item \textbf{Ripetere}
\end{enumerate}

\section{Tipi di test}
\begin{itemize}
    \item \textbf{Test di unità}\\
    Si concentrano sulle singole unità di codice, ovvero le singole classi e i metodi
    \item \textbf{Test di integrazione}\\
    Verificano che i vari moduli, classi o componenti comunichino e funzionino correttamente quando vengono combinati
    \item \textbf{Test di sistema}\\
    Valutano il prodotto software nella sua interezza per verificare che rispeddi le specifiche architettoniche e globali
    \item \textbf{Test di accettazione}\\
    Dimostra che il sistema soddisfa effettivamente i requisiti e i bisogni dell'utente finale o del committente
\end{itemize}


\chapter{Code Smell}
\section{Cosa sono?}
I Code Smell sono indicatori di problemi strutturali profondi.

\section{Long Method}
\begin{table}[h]
    \centering
    \renewcommand{\arraystretch}{1.2} % Migliora l'altezza delle righe
    \begin{tabularx}{\textwidth}{l X}
        \toprule
        \textbf{Nome} & Long Method \\
        \midrule
        \textbf{Problema} & Un metodo svolge troppe azioni e supera le 10 righe di codice, diventando difficile da capire e riutilizzare. \\
        \midrule
        \textbf{Soluzione} & 
        \begin{itemize}
            \item \textbf{Extract Method:} Si isolano i blocchi logici e si creano dei nuovi metodi che li contengono.
            \item \textbf{Replace Temp with query:} Si sposta l'espressione in un metodo separato che ne restituisce il risultato.
            \item \textbf{Introduce Parameter Object:} Si raggruppano i parametri ripetuti in un oggetto dedicato.
        \end{itemize}
        \\ % Questo backslash è FONDAMENTALE per chiudere la riga della tabella
        \bottomrule
    \end{tabularx}
\end{table}

\section{Large Class}
\begin{table}[h]
    \centering
    \renewcommand{\arraystretch}{1.2} % Migliora l'altezza delle righe
    \begin{tabularx}{\textwidth}{l X}
        \toprule
        \textbf{Nome} & Large Class 
        \\
        \midrule
        \textbf{Problema} &
        Una classe contiene molti campi/metodi/linee di codice
        \\
        \midrule
        \textbf{Soluzione} & 
        \begin{itemize}
            \item \textbf{Extract Class:} Prende un sottoinsieme delle variabili di istanza e dei metodi e crea una nuova classe con essi.
            
            Questo rende la classe iniziale più breve
            \item \textbf{Move Method:} Sposta uno o più metodi in altre classi
            \item \textbf{Extract Subclass:} Crea una sottoclasse per una parte del comportamento della classe, se questo è usato solo in casi
        \end{itemize}
        \\ % Questo backslash è FONDAMENTALE per chiudere la riga della tabella
        \bottomrule
    \end{tabularx}
\end{table}

\section{Duplicated Code}
\begin{table}[h]
    \centering
    \renewcommand{\arraystretch}{1.2} % Migliora l'altezza delle righe
    \begin{tabularx}{\textwidth}{l X}
        \toprule
        \textbf{Nome} & Duplicated Code \\
        \midrule
        \textbf{Problema} & Lo stesso codice o molto simile appare in più istanze \\
        \midrule
        \textbf{Soluzione} & 
        Se lo stesso codice si trova in 2 o più metodi della stessa classe usare Extract Method e posizionare le chiamate per il nuovo metodo in entrami i posti
        \\ % Questo backslash è FONDAMENTALE per chiudere la riga della tabella
        \bottomrule
    \end{tabularx}
\end{table}

\newpage
\section{Feature Envy}
\begin{table}[h]
    \centering
    \renewcommand{\arraystretch}{1.2} % Migliora l'altezza delle righe
    \begin{tabularx}{\textwidth}{l X}
        \toprule
        \textbf{Nome} & Feature Envy \\
        \midrule
        \textbf{Problema} & Un metodo accede ai dati di un altro oggetto più che ai propri \\
        \midrule
        \textbf{Soluzione} & 
        \begin{itemize}
            \item Se un metodo deve essere chiaramente spostato in un altro luogo, utilizzare Move Method
            \item Se solo una parte di un metodo accede ai dati di un altro oggetto, usare Extract Method
            \item Se un metodo utilizza funzioni di diverse altri classi, determinare innanzitutto quale classe contiene la maggior parte dei dati utilizzaati.

            Poi collocare il metodo in questa classe insieme agli altri dati.
            \item Usare Extract MEthod per dividere il metodo in diverse parti che possono essere collocate in diversi posti in diverse classi
        \end{itemize}
        \\ % Questo backslash è FONDAMENTALE per chiudere la riga della tabella
        \bottomrule
    \end{tabularx}
\end{table}

\section{Switch Statement}
\begin{table}[h]
    \centering
    \renewcommand{\arraystretch}{1.2} % Migliora l'altezza delle righe
    \begin{tabularx}{\textwidth}{l X}
        \toprule
        \textbf{Nome} & Switch Statement \\
        \midrule
        \textbf{Problema} & Si ha un complesso operatore di Switch o di una sequenza di istruzioni if \\
        \midrule
        \textbf{Soluzione} & 
        \begin{itemize}[leftmargin=*, nosep, after=\vspace{0.5em}]
           \item \textbf{Replace Conditional with Polymorphism:} Creare sottoclassi corrispondenti ai rami della condizione.

           In esse, creare un metodo condiviso e spostare il codice dal ramo corrispondente della condizione di esso.

           Quindi sostituire il condizionale con la chiamata al metodo corrispondente.

           Il risultato è che l'implementazione corretta sarà ottenuta tramite il polimorfismo, a seconda della classe dell'oggetto
           \item Per isolare lo "switch" e metterlo nella sottoclasse giusta, potrebbe essere necessario Extract Method e poi Move Method
        \end{itemize}
        \\ % Questo backslash è FONDAMENTALE per chiudere la riga della tabella
        \bottomrule
    \end{tabularx}
\end{table}

\section{Data Class}
\begin{table}[h]
    \centering
    \renewcommand{\arraystretch}{1.2} % Migliora l'altezza delle righe
    \begin{tabularx}{\textwidth}{l X}
        \toprule
        \textbf{Nome} & Data Class \\
        \midrule
        \textbf{Problema} &
        Una classe che ha solo variabili di classe, metodi/proprietà, getter/setter e nient'altro.
        
        Sta solo agendo come contenitore di dati
        \\
        \midrule
        \textbf{Soluzione} & 
        Esaminare i metodi che utilizzano i dati nella data class e usare Move Method o Extract Method per migrare questa funzionalità nella data class
        \\ % Questo backslash è FONDAMENTALE per chiudere la riga della tabella
        \bottomrule
    \end{tabularx}
\end{table}
 \newpage
\section{Long Parameter List}
\begin{table}[h]
    \centering
    \renewcommand{\arraystretch}{1.2} % Migliora l'altezza delle righe
    \begin{tabularx}{\textwidth}{l X}
        \toprule
        \textbf{Nome} & Long Parameter List \\
        \midrule
        \textbf{Problema} & Molti parametri passati in un metodo \\
        \midrule
        \textbf{Soluzione} & 
        \begin{itemize}[leftmargin=*, nosep, after=\vspace{0.5em}]
           \item Se alcuni degli argomenti sono solo risultati di chiamate di metodo di un altro oggetto usare Replace Parameter with Method Call
           \item Invece di passare un gruppo di dati ricevuti da un altro oggetto come parametri, passare l'oggetto stesso al metodo, utilizzando Preserve Whole Object
           \item Se ci sono più elementi dati non correlati tra loro, a volte è possibile unirli in un unico oggetto parametro tramite Introduce Parameter Object
        \end{itemize}
        \\ % Questo backslash è FONDAMENTALE per chiudere la riga della tabella
        \bottomrule
    \end{tabularx}
\end{table}

\section{Shotgun Surgery}
\begin{table}[h]
    \centering
    \renewcommand{\arraystretch}{1.2} % Migliora l'altezza delle righe
    \begin{tabularx}{\textwidth}{l X}
        \toprule
        \textbf{Nome} & Shotgun Surgery \\
        \midrule
        \textbf{Problema} & Per apportare eventuali modifiche è necessario apportare molte piccole modifiche a molte classi diverse \\
        \midrule
        \textbf{Soluzione} & 
        \begin{itemize}[leftmargin=*, nosep, after=\vspace{0.5em}]
            \item Usare il metodo Move Method e Move Field per spostare i comportamenti di classe esistenti in una singola classe.

            Se non c'è classe appropriata per questo, crearne una nuova
            \item Se lo spostamento del codice nella stessa clases lascia le classi originali quasi vuote, provare a sbarazzarsi di queste classi ormai ridondati tramite Inline Class
        \end{itemize}
        \\ % Questo backslash è FONDAMENTALE per chiudere la riga della tabella
        \bottomrule
    \end{tabularx}
\end{table}

\section{Comment}
\begin{table}[h]
    \centering
    \renewcommand{\arraystretch}{1.2} % Migliora l'altezza delle righe
    \begin{tabularx}{\textwidth}{l X}
        \toprule
        \textbf{Nome} & Commentthod \\
        \midrule
        \textbf{Problema} & Un metodo è ricco di commenti esplicativi\\
        \midrule
        \textbf{Soluzione} & 
        \begin{itemize}[leftmargin=*, nosep, after=\vspace{0.5em}]
            \item Se un commento è inteso a spiegare un'interrogaione complessa, l'espressione deve essere suddivisa in sottoespressioni comprensibili utilizzando Extract Variable
            \item Se un commento spiega una sezione di codice, questa sezione può essere trasformata in un medoto separato tramite Extract MEthod
            \item Se un metodo è già stato estratto, ma sono ancora necessari commenti per spiegare cosa fa il metodo, date al metodo un nome autoesplicativo con Rename Method
            \item Se avete bisgono di affermare delle regole su uno stato è necessario per il funzionamento del sistemare, usate Introduce Assertion
        \end{itemize}
        \\ % Questo backslash è FONDAMENTALE per chiudere la riga della tabella
        \bottomrule
    \end{tabularx}
\end{table}

\newpage
\section{Refused Bequest}
\begin{table}[h]
    \centering
    \renewcommand{\arraystretch}{1.2} % Migliora l'altezza delle righe
    \begin{tabularx}{\textwidth}{l X}
        \toprule
        \textbf{Nome} & Refused Bequest \\
        \midrule
        \textbf{Problema} & Una sottoclasse utilizza solo alcuni dei metodi e delle proprietà ereditati dai suoi genitori, la gerarchia non è corretta \\
        \midrule
        \textbf{Soluzione} & 
        \begin{itemize}[leftmargin=*, nosep, after=\vspace{0.5em}]
            \item \textbf{Replace Integration with Delegation:} Creare un campo e inserire un oggetto della superclasse, delegare i metodi all'oggetto della superclasse e rimuore l'ereditarietà
            \item \textbf{Extract Superclass:} Creare una superclasse condivisa e trasferirvi tutti i campi e metodi comuni
        \end{itemize}
        \\ % Questo backslash è FONDAMENTALE per chiudere la riga della tabella
        \bottomrule
    \end{tabularx}
\end{table}